% Penggerak mandiri untuk bab Aritmetisasi dalam pelokalan Bahasa Indonesia.

\documentclass[openany]{memoir}

% `cleveref` Open Logic belum memiliki opsi Indonesian. Babel tetap memuat
% Indonesian; `\ollanguage` bernilai english hanya selama cleveref dimuat.
\PassOptionsToPackage{english,indonesian}{babel}

\newcommand*{\olpath}{../..}
\newcommand*{\ollangid}{id}
\newcommand*{\ollanguage}{english}

\usepackage{mathpazo}
\usepackage{helvet}
\usepackage{courier}
\linespread{1.05}

\input{\olpath/sty/open-logic.sty}
\renewcommand*{\ollanguage}{indonesian}
\input{\olpath/sty/open-logic-defer.sty}

\problemsperchapter
\let\cleardoublepage\clearpage
\pagestyle{plain}

\begin{document}

\selectlanguage{indonesian}
% Referensi ke dua bagian Historis yang muncul kemudian dalam pembaca lengkap
% dideklarasikan sebagai metadata saja. Sumber Inggrisnya tidak dimasukkan ke
% edisi id-ID ini, sehingga tidak ada fallback prosa Inggris.
\makeatletter
\hypertarget{ol-id-forward-history}{}
\protected@write\@auxout{}{\string\newlabel{his:set:limits:sec}{{H.1}{0}{Definisi Limit yang Ketat}{ol-id-forward-history}{}}}
\protected@write\@auxout{}{\string\newlabel{his:set:limits:sec@cref}{{[section][1][]H.1}{[1][0][]0}{}{}{}}}
\protected@write\@auxout{}{\string\newlabel{his:set:mythology:sec}{{H.2}{0}{Lebih Banyak Mitos daripada Sejarah?}{ol-id-forward-history}{}}}
\protected@write\@auxout{}{\string\newlabel{his:set:mythology:sec@cref}{{[section][2][]H.2}{[1][0][]0}{}{}{}}}
\makeatother
% Bab-bab sebelumnya disertakan agar semua rujukan silang dalam bab
% Aritmetisasi terurai dari sumber id-ID tanpa fallback bahasa Inggris.
\olimport*[sets-functions-relations/sets]{sets}
\olimport*[sets-functions-relations/relations]{relations-complete}
\olimport*[sets-functions-relations/functions]{functions}
\olimport*[sets-functions-relations/size-of-sets]{size-of-sets-complete}
\olimport*[sets-functions-relations/arithmetization]{arithmetization}

\bibliographystyle{\olpath/bib/natbib-oup}
\bibliography{\olpath/bib/open-logic}

\end{document}
